\documentclass[../main.tex]{subfiles}

\begin{document}

\section{意味論と構文論}

論理学では，言語表現そのものを扱う構文論と，言語表現が何を指すかを扱う意味論を区別する。構文論は記号の形や結合規則を扱い，意味論はそれらの記号が世界のどの対象，集合，真理値を指示するかを扱う。

\begin{center}
\begin{tikzpicture}[>=Stealth, node distance=2cm]
  \node[draw, ellipse, minimum width=3cm, minimum height=5cm] (sem) at (0,0) {};
  \node at (0,1.3) {実在};
  \node at (0,0.7) {世界};
  \node at (0,0.1) {個体};
  \node at (0,-0.8) {集合};
  \node at (0,-1.8) {真理値};

  \node[draw, ellipse, minimum width=3cm, minimum height=4cm] (syn) at (6,0) {};
  \node at (6,1.1) {言語};
  \node at (6,0.5) {項};
  \node at (6,-0.2) {述語関数};
  \node at (6,-1.1) {命題};

  \draw[->] (4.8,0.5) -- (1.3,0.1);
  \draw[->] (4.8,-0.2) -- (1.3,-0.8);
  \draw[->] (4.8,-1.1) -- (1.3,-1.8);

  \node at (3.1,-0.45) {指示 \((\mathrm{refer})\)};
  \node at (0,-3) {semantics};
  \node at (6,-3) {syntax};
\end{tikzpicture}
\end{center}

意味論とは，言葉がどのような形で何を指しているかを明らかにする理論である。上の図では，項は個体を，述語関数は集合を，命題は真理値を指示するものとして対応づけられている。

\subsection{真理関数的文脈の例}

命題 \(P,Q\) から新しい命題を作る文脈として，次のものがある。
\[
  P \text{ または } Q,
  \qquad
  P \text{ でない},
  \qquad
  P \text{ かつ } Q
\]
これらの文脈は，入力される命題の真理値だけによって出力の真理値が決まる。

\begin{center}
\begin{tabular}{ccc}
\toprule
\(P\) & \(Q\) & \(P \text{ または } Q\) \\
\midrule
1 & 1 & 1 \\
1 & 0 & 1 \\
0 & 1 & 1 \\
0 & 0 & 0 \\
\bottomrule
\end{tabular}
\qquad
\begin{tabular}{cc}
\toprule
\(P\) & \(P \text{ でない}\) \\
\midrule
1 & 0 \\
0 & 1 \\
\bottomrule
\end{tabular}
\qquad
\begin{tabular}{ccc}
\toprule
\(P\) & \(Q\) & \(P \text{ かつ } Q\) \\
\midrule
1 & 1 & 1 \\
1 & 0 & 0 \\
0 & 1 & 0 \\
0 & 0 & 0 \\
\bottomrule
\end{tabular}
\end{center}

\section{関係}

関係とは，複数の対象のあいだに成り立つ結びつきである。特に2つの対象 \(x,y\) のあいだに関係 \(R\) が成り立つことを
\[
  xRy
\]
と表し，これを2項関係という。

\subsection{関係の法則}

関係 \(R\) の基本的な性質として，反射律，対称律，推移律がある。

\paragraph{反射律}
任意の \(x \in D\) に関して，
\[
  xRx
\]
が成り立つとき，\(R\) は反射的であるという。

\paragraph{対称律}
任意の \(x,y \in D\) に関して，
\[
  xRy \quad \Rightarrow \quad yRx
\]
が成り立つとき，\(R\) は対称的であるという。講義ノートではこの対応が \(xRy = yRx\) と記されていた。

\paragraph{推移律}
任意の \(x,y,z \in D\) に関して，
\[
  xRy \ \text{かつ}\ yRz
  \quad \Rightarrow \quad
  xRz
\]
が成り立つとき，\(R\) は推移的であるという。

\subsection{代表的な関係の性質}

代表的な関係について，反射律・対称律・推移律が成り立つかどうかを整理すると，次のようになる。ただし，元資料では列見出しと一部の記号が不鮮明であったため，判読できる情報を残している。

\begin{center}
\begin{tabular}{c|ccccccc}
\toprule
 & \(\in\) & \(<\) & \(\neq\) & \(\leq\) & \(\equiv\) & \(=\) & 判読不明 \\
\midrule
反射律 & \(\times\) & \(\times\) & \(\times\) & \(\circ\) & \(\times\) & \(\circ\) & \(\circ\) \\
対称律 & \(\times\) & \(\times\) & \(\circ\) & \(\times\) & \(\circ\) & \(\times\) & \(\circ\) \\
推移律 & \(\times\) & \(\circ\) & \(\times\) & \(\times\) & \(\circ\) & \(\circ\) & \(\times\) \\
\bottomrule
\end{tabular}
\end{center}

\begin{sidebox}[TBred]{判読注意}
  写真では列見出しが一部不鮮明であり，\(\in,<,\neq,\leq,\equiv,=\) 付近に見える。元の表には各行7個の判定が残っているため，最後の列は判読不明として保持した。
\end{sidebox}

\subsection{束論への接続}

関係を矢印で表すと，対象のあいだの順序構造を図として見ることができる。次のような有向図は，後に扱う束論への導入になる。

\begin{center}
\begin{tikzpicture}[>=Stealth, node distance=1.2cm]
  \node (a0) at (0,0) {\(a_0\)};
  \node (a1) at (1.5,0) {\(a_1\)};
  \node (a2) at (3,1) {\(a_2\)};
  \node (a3) at (3,-1) {\(a_3\)};
  \node (a7) at (5,1.8) {\(a_7\)};
  \node (a8) at (5,1.1) {\(a_8\)};
  \node (a9) at (5,0.4) {\(a_9\)};
  \node (a4) at (5,-0.5) {\(a_4\)};
  \node (a5) at (5,-1.3) {\(a_5\)};
  \node (a6) at (7,-0.8) {\(a_6\)};
  \node (a10) at (7,1.1) {\(a_{10}\)};
  \node (a11) at (8.5,0.1) {\(a_{11}\)};

  \draw[->] (a0) -- (a1);
  \draw[->] (a1) -- (a2);
  \draw[->] (a1) -- (a3);
  \draw[->] (a2) -- (a7);
  \draw[->] (a2) -- (a8);
  \draw[->] (a2) -- (a9);
  \draw[->] (a8) -- (a10);
  \draw[->] (a10) -- (a11);
  \draw[->] (a3) -- (a4);
  \draw[->] (a3) -- (a5);
  \draw[->] (a4) -- (a6);
  \draw[->] (a5) -- (a6);
  \draw[->] (a6) -- (a11);

  \node[text=TBred] at (8,-1.2) {束論};
  \node[text=TBred] at (8,-1.65) {lattice};
\end{tikzpicture}
\end{center}

\section{順序構造}

\subsection{半順序構造}

集合 \(M\) とその上の関係 \(\leq\) の組
\[
  \mathcal{S}
  =
  \langle M,\leq\rangle
\]
を考える。この構造が反射律と推移律を満たすとき，半順序構造という。英語では partial ordered structure と呼ばれ，略して p.o.s. と書かれる。
\[
  (\mathrm{p.o.s.};\ \mathrm{partial\ ordered\ structure})
\]

\paragraph{反射律}
任意の \(x \in M\) について，
\[
  x \leq x
\]
が成り立つ。

\paragraph{推移律}
任意の \(x,y,z \in M\) について，
\[
  x \leq y
  \quad \text{かつ} \quad
  y \leq z
  \quad \Rightarrow \quad
  x \leq z
\]
が成り立つ。

\subsection{全順序構造}

全順序構造とは，半順序構造の条件に加えて，任意の2つの元が比較可能であるという条件を満たす構造である。
\[
  \text{全順序}
  \quad
  (\mathrm{t.o.s.};\ \mathrm{total\ ordered\ structure})
\]

\paragraph{全順序律}
任意の \(x,y \in M\) について，
\[
  x \leq y
  \quad \text{または} \quad
  y \leq x
\]
が成り立つ。

\subsection{等号の導入}

順序関係 \(\leq\) を用いて等号を導入することができる。任意の \(x,y \in M\) について，
\[
  x \leq y
  \quad \text{かつ} \quad
  y \leq x
  \quad \Leftrightarrow \quad
  x = y
\]
と定める。このとき，等号 \(=\) は \(\leq\) から導かれる関係であり，反射律・対称律・推移律をすべて満たす関係であることが分かる。

\end{document}
